\manualtablecaption{System-Register Builtins}
\begingroup\footnotesize
\setlength{\tabcolsep}{2pt}
\begin{longtable}{@{}p{1.4in}p{1.3in}p{0.7in}p{2.2in}@{}}
\toprule
\textbf{Name} & \textbf{C interface} & \textbf{Lowering} & \textbf{Availability, constraint, and effect}\\
\midrule
\endhead
\texttt{write\_status} & \texttt{void(uint16\_t)} & \texttt{WRSTATUS} & Supervisor; reserved and privilege-controlled bits follow the ISA; compiler-visible side effect.\\
\texttt{read\_control\_register} & \texttt{uint64\_t(selector)} & \texttt{RDCR} & Supervisor; selector is a constant in 0 through 65535 naming a defined register.\\
\texttt{write\_control\_register} & \texttt{void(selector,u64)} & \texttt{WRCR} & Supervisor; defined constant selector required; writes the register and clobbers memory.\\
\texttt{read\_segment\_register} & \texttt{uint64\_t(selector)} & \texttt{RDSEG} & Unprivileged; selector is a compile-time \texttt{\_\_bedrock\_segment\_register\_t}.\\
\texttt{read\_code\_segment} & \texttt{uint64\_t(void)} & \texttt{RDSEG CS} & Unprivileged; reads the current CS image through the dedicated source-only encoding.\\
\texttt{write\_segment\_register} & \texttt{void(selector,image)} & \texttt{WRSEG} & Unprivileged; validates the SREG image, changes address interpretation, and fully clobbers memory.\\
\bottomrule
\end{longtable}
\endgroup
